% Generated draft: requires mathematical and source review.
\documentclass[11pt]{article}
\usepackage[margin=25mm]{geometry}
\usepackage{fontspec}
\setmainfont{Latin Modern Roman}
\usepackage{amsmath,amssymb,mathtools,unicode-math}
\setmathfont{Latin Modern Math}
\usepackage{enumitem,needspace,xurl,hyperref}
\setlength{\parindent}{0pt}
\setlength{\parskip}{5pt}
\setlength{\emergencystretch}{3em}
\newcommand{\N}{\mathbb{N}}
\newcommand{\Z}{\mathbb{Z}}
\newcommand{\Q}{\mathbb{Q}}
\newcommand{\R}{\mathbb{R}}
\newcommand{\C}{\mathbb{C}}
\newcommand{\F}{\mathbb{F}}
\newcommand{\A}{\mathbb{A}}
\newcommand{\PP}{\mathbb{P}}
\newcommand{\E}{\mathbb{E}}
\newcommand{\eps}{\varepsilon}
\newcommand{\dd}{\,\mathrm{d}}
\begin{document}
\section*{UnsolvedMath ID: 30004483}
\textbf{Source problem code:} catalogue entry
\subsection{Definitions and mathematical statement}
Let \(T\) be an area-minimizing two-dimensional current with analytic prescribed boundary curve \(\Gamma\subset\R^{2+n}\), as in the supplied source. Write \(\partial T\) for the boundary current and \(\operatorname{spt}T\) for the support of \(T\). The source does not specify the range of \(n\), the coefficient group or admissible class of currents, the precise minimization condition, or the orientation and multiplicity of the prescribed boundary; consequently, no particular identity between \(\partial T\) and an integration current over \(\Gamma\) is imposed here.

Let \(\mathcal B_{\Gamma}(T)\subseteq\Gamma\cap\operatorname{spt}T\) denote the set called the \emph{branched singular points at the boundary} in the source. This notation names the intended set; the supplied record does not give a criterion defining membership in it. For \(p\in\R^{2+n}\) and \(r>0\), put
\[
B_r(p)=\{x\in\R^{2+n}:|x-p|<r\}.
\]
The isolation part of the proposed assertion is
\[
\forall p\in\mathcal B_{\Gamma}(T)\quad
\exists r_p>0\quad
\mathcal B_{\Gamma}(T)\cap B_{r_p}(p)=\{p\}.
\]
The complete source question asks whether this isolation assertion holds for every such \(T\), and whether \(T\) also has \emph{finite topology locally at every} \(p\in\mathcal B_{\Gamma}(T)\).

The latter phrase is retained as an unresolved qualification: the record does not specify the topological object associated with the current, the meaning of finite topology, or the required neighborhood condition. It therefore does not support replacing this phrase by a particular finite-genus, finite-component, or finite-type condition. The precise definitions of boundary branching, analytic boundary curve, admissible current, and local finite topology require confirmation from the source before this becomes a fully specified conjecture.
\subsection{Short English statement}
For an area-minimizing two-dimensional current with analytic prescribed boundary, must every boundary branch singularity be isolated, and must the current have finite topology near each such point? The supplied record labels this question open, but several technical conventions needed for a fully precise statement are unspecified.
\subsection{Sources}
\begin{itemize}
\item Local registry: supplied UnsolvedMath record, permanent ID \(30004483\), title \emph{Boundary Branch Singularities of Area-Minimizing Currents}, source problem code \texttt{catalogue entry}, in the Oberwolfach Reports Open Problems collection. No individual proposer is supplied; the proposed year is recorded as 2020.
\item Supplied report citation: \emph{Calculus of Variations} (2021), Oberwolfach Reports, pp.\ 1139--1196, DOI \texttt{10.4171/owr/2020/22}. \url{https://doi.org/10.4171/owr/2020/22}. The supplied background describes these reports as documenting mathematical workshops and their open-problem discussions.
\item Supplied literature source: OWR Calculus of Variations report. \url{https://ems.press/content/serial-article-files/46860}.
\item Supplied literature source: De Lellis--De Philippis--Hirsch, \emph{Nonclassical minimizing surfaces with smooth boundary}. \url{https://arxiv.org/abs/1906.09488}.
\item Supplied literature assessment, checked \texttt{2026-08-22}: the record labels the problem open and states that the linear Dirichlet-energy model has discrete singular set results for analytic boundary curves, while isolation and finite local topology for general area-minimizing currents remain presented as a conjecture. This is supplied dated curation derived from the OWR material, not an independently verified literature assessment. The record requests independent verification of time-sensitive status and cited claims.
\end{itemize}
\end{document}
